\documentclass[12pt]{article}
\usepackage[utf8]{inputenc}
\usepackage[T1]{fontenc}
\usepackage{hyperref}
\title{\textbar Xam Language Resources for K-12 Curriculum Development}
\author{Compiled by Claude Sonnet 5.5 (claude-sonnet-5-5) for Andr\'as Kornai}
\date{}
\begin{document}
\maketitle

\section{Introduction}
\textbar Xam (ISO 639-3 xam, Glottocode xamm1241) is an extinct Tuu language of South Africa, and it has no known native speakers. It is documented almost entirely in the Bleek and Lloyd Collection. Wilhelm Bleek and Lucy Lloyd recorded it from 1870 onward from \textbar xam speakers, including \textbar\textbar kabbo, Dia!kwain and \textbar han$\neq$kass'o. A heritage and revitalization curriculum would therefore be built from archival material. No formal school curriculum was found in these searches.

\section{Primary Educational Resources}
\begin{itemize}
\item Pippa Skotnes, \emph{When the World Was}, a book for children and adults drawn from the stories told to Bleek and Lloyd: \url{https://humanities.uct.ac.za/cca/articles/2025-06-05-when-world-was-book-children-and-adults}
\item ``High Sky Stories'' report. It describes a Children's Book Network workshop in Clanwilliam where storytellers told San stories from the Bleek and Lloyd collection: \url{https://www.childrensbook.co.za/2013/05/20/high-sky-stories-short-report/}
\item Lesley Beake's picture book \emph{Little Lion} (2011) also draws on Bleek and Lloyd material, but no URL for it appeared in the results.
\end{itemize}

\section{Linguistic and Reference Materials}
\begin{itemize}
\item Wikipedia article on the language: \url{https://en.wikipedia.org/wiki/%C7%80Xam_language}
\item About the Digital Bleek and Lloyd: \url{https://bal2.cs.uct.ac.za/about.html}
\item Digital Bleek and Lloyd manuscript dictionary (\textbar Xam index): \url{https://digitalbleeklloyd.uct.ac.za/metadata/dictionaries/Xam/index.html}
\item The archive holds 15,325 notebook spreads and 85,821 dictionary cards, plus over 500 drawings.
\item \emph{Specimens of Bushman Folklore}: \url{https://en.wikipedia.org/wiki/Specimens_of_Bushman_Folklore}
\item Dorothea Bleek, \emph{Customs and Beliefs of the \textbar Xam} (Wits University Press, 2nd ed.\ 2022). The only URL seen was a library catalogue record, so it is not included in the JSON.
\end{itemize}

\section{Religious and Community Resources}
No religious texts or community organisations specific to \textbar Xam were found in these searches. The stories in the collection are traditional narratives and beliefs.

\section{Digital and Audio Resources}
The Digital Bleek and Lloyd site, run by the Centre for Curating the Archive at the University of Cape Town, provides the digitised notebooks, dictionary cards and stories: \url{https://lloydbleekcollection.cs.uct.ac.za/stories/1127/index.html}. No audio recordings were found, which is expected for a language whose last speakers died long ago.

\section{Curriculum Development Context}
No official heritage curriculum was found. A K-12 programme would have to draw on the stories and vocabulary in the archive, and on the existing children's books and community storytelling projects above. The Department of Basic Education's CAPS documents and the UCT project are the next places to check.

\end{document}